\section{Key Distributions in Practice}
\label{sec:distributions}

The theoretical bounds of Section~\ref{sec:bounds} hold for any key set.
In practice, keys in storage systems follow distributions shaped
by the application domain.

\subsection{Common Industrial Distributions}
\label{subsec:common-industrial-distributions}

Keys in LSM-tree storage systems typically exhibit one or more of the following patterns:

\paragraph{Sequential / near-sequential.}
Auto-incrementing primary keys, monotonic timestamps, log sequence numbers. Keys are
approximately evenly spaced with small gaps.
The key space is dense – the spread is much smaller than the universe.
Among the Search on Sorted Data (SOSD)~\citep{kipf2019sosd} datasets, Facebook (user IDs), Books (popularity keys), and Wiki
(edit timestamps) all fall into this category.

\paragraph{Clustered.}
Keys concentrate in a few dense regions separated by large gaps.
Event logs keyed by user ID prefix + timestamp cluster around active users; inactive users leave gaps.
Geospatial keys such as S2 cell IDs~\citep{s2geometry} cluster around populated areas, with sparse rural regions.
OpenStreetMap (OSM) cell IDs are a representative example: they span nearly the full 64-bit space,
yet most keys concentrate in a small fraction of it.

\paragraph{Uniform-like.}
UUIDs, cryptographic hashes, randomly generated tokens.
Keys are spread across the full key space with large, roughly equal gaps.

\subsection{SOSD Benchmark Datasets}
\label{sec:sosd}

For empirical evaluation, we use real-world datasets from the SOSD (Search on Sorted
Data) benchmark~\citep{kipf2019sosd}:

\begin{table}[H]
    \centering
    \begin{tabular}{@{}llp{0.68\textwidth}@{}}
        \toprule
        Dataset & $n$ (full) & Character \\
        \midrule
        \textbf{Facebook} & 200M &
        User IDs from Facebook's social graph (uint64). \\
        \textbf{Books} & 200M &
        Amazon book popularity keys (uint32). \\
        \textbf{Wiki} & 200M &
        Unix timestamps of Wikipedia article edits (uint64). \\
        \textbf{OSM} & 800M &
        Hierarchical cell IDs covering global OpenStreetMap data (uint64). \\
        \bottomrule
    \end{tabular}
    \caption{SOSD datasets used in our evaluation.
    Benchmarks in Chapter~\ref{chap:evaluation} use
    subsets of these datasets (e.g.\ $n = 2^{24}$)}
    \label{tab:sosd}
\end{table}

\begin{remark}
    All datasets are sorted and deduplicated on load.
    The Wiki dataset contains many repeated timestamps, so its effective $n$
    is smaller than the requested raw count.
\end{remark}

\begin{table}[H]
    \centering
    \begin{tabular*}{\textwidth}{@{\extracolsep{\fill}}lrcrrrrrr@{}}
        \toprule
        Dataset  & $n$                   & $U$ [bits]    & P10               & P50               & P90               & max               & mean \\
        \midrule
        Facebook & ${\approx}200\text{M}$ & 37           & 2                 & 28                & 693               & $5.7\times10^{5}$ & 387 \\
        Wiki     & ${\approx}90\text{M}$  & 29           & 1                 & 1                 & 3                 & $1.7\times10^{5}$ & 3 \\
        Books    & ${\approx}200\text{M}$ & 31           & 1                 & 2                 & 9                 & $2.0\times10^{4}$ & 5 \\
        OSM      & 800M                   & 64           & $1.6\times10^{6}$ & $4.5\times10^{7}$ & $1.7\times10^{9}$ & $2.3\times10^{17}$ & $1.7\times10^{10}$ \\
        \bottomrule
    \end{tabular*}
    \caption{Consecutive-key gap statistics for SOSD datasets after sort and deduplication.
        $n$: effective key count.
        $U$: key range in bits ($\lfloor\log_2(\max-\min)\rfloor + 1$).
        Gaps are differences between adjacent sorted keys.}
    \label{tab:sosd-gaps}
\end{table}

\begin{figure}[H]
    \centering
    \begin{tabular}{@{}cc@{}}
        \includegraphics[width=0.48\textwidth]{hist_sosd_fb} &
        \includegraphics[width=0.48\textwidth]{hist_sosd_books} \\
        \includegraphics[width=0.48\textwidth]{hist_sosd_wiki} &
        \includegraphics[width=0.48\textwidth]{hist_sosd_osm} \\
    \end{tabular}
    \caption{Key density histograms for SOSD datasets (1000 bins, log-scale Y-axis).
        Relative density is the fraction of all keys falling in each bin,
        normalized so that a perfectly uniform distribution gives~1 everywhere.}
    \label{fig:sosd-histograms}
\end{figure}

\subsection{Synthetic Distributions}
\label{subsec:synthetic-distributions}

To isolate the effect of specific distribution properties, we also generate synthetic
key sets:

\begin{itemize}
    \item \textbf{Uniform:} keys drawn uniformly from $[0, 2^{60})$.
    Average gap $\approx 2^{60}/n$.
    No clustering.
    \item \textbf{Clustered:} a mixture of Gaussian clusters with controlled parameters
    (number of clusters, intra-cluster spread, inter-cluster gap).
    \colorbox{orange}{TODO: state the concrete parameter values used in the benchmarks
        (number of clusters, $\sigma$, inter-cluster gap) --- see
        \texttt{Thesis/testutils/distribution.go:GenerateClusterDistribution}.}
\end{itemize}

Together, the SOSD and synthetic datasets span the spectrum from near-sequential
(Facebook, Books, Wiki) through clustered (OSM) to structureless (uniform
random) --- allowing us to evaluate each filter design across the range of distributions
it may encounter in practice.

\subsection{An assumption on queries}
\label{subsec:assumption}

We assume that queries target regions of the key space where stored keys are dense,
rather than arbitrary empty regions.

\begin{assumption}
    Queries reflect the application state and therefore concentrate on
    regions that contain data.
\end{assumption}

Consequently, FPR in sparse regions of the key space matters less --- dense regions
dominate the effective FPR experienced by the application.

The SODA hash (Section~\ref{sec:soda}) ignores this assumption entirely --- it provides
worst-case guarantees for any distribution. The truncation hash
(Section~\ref{sec:truncation}) performs well only under uniformity. The hybrid approaches
(Chapter~\ref{chap:hybrid}) explicitly exploit it by isolating dense regions.

Section~\ref{sec:query-distribution} operationalizes this assumption
in the query workload.

\section{Query generation}\label{sec:query-distribution}

Queries are drawn from a three-bucket mixture that reflects
the assumption of Section~\ref{subsec:assumption}:

\begin{itemize}
    \item \textbf{50\% near-key.}
    A random stored key $k$ is chosen; the query start is drawn
    uniformly from $[k - 5L,\; k + 5L)$.
    Models application lookups that miss by a small margin.
    \item \textbf{30\% in-gap.}
    A random gap between two adjacent keys is selected uniformly;
    the query is placed entirely inside it.
    Models misses in locally dense regions.
    \item \textbf{20\% uniform.}
    A random start is drawn uniformly from $[\min k, \max k]$.
    Covers sparse regions to verify FPR stays reasonable there too.
\end{itemize}

The near-key and in-gap buckets reflect the dominant workload;
the uniform bucket acts as a safety check on sparse regions.
The headline \emph{gap-heavy} workload of
Section~\ref{sec:eval-tables} reweighs the mixture to
$0\,/\,70\,/\,30$ (no near-key, 70\% in-gap, 30\% uniform) so that no
candidate query is guaranteed empty just by its construction.

The two metrics this thesis cares about --- FPR and query latency
--- want different things from the same query stream, so we draw
two query sets from the mixture above and let truncation decide
which.

\subsection{Empty queries for False Positive Ration Measurement}
\label{subsec:queries-empty}

For FPR estimation every measured query must be empty: a positive
answer on a non-empty query is a true positive, not a false one,
and would deflate the FPR estimate. We enforce emptiness by
\emph{truncation}: when a candidate query $[a, a+L)$ would contain
a stored key $k^*$, it is shortened to $[a, k^* - 1]$ rather than
discarded and resampled. Truncation preserves queries that start
inside a cluster and end between two closely spaced keys, which
would otherwise be rejected every time --- biasing the workload
towards cluster boundaries and missing the harder case of a query
landing deep inside a dense region. The effective range length
may therefore be smaller than $L$.

\subsection{Combined queries (for latency)}
\label{subsec:queries-combined}

For query-latency measurements we want the workload to look like
the actual LSM read pattern: a fraction of the queries land on a
stored key and the filter answers ``not empty'' truthfully, not as
a false positive.

The combined-queries generator draws from the same three-bucket
mixture and skips the truncation step, so a query may contain one
or more stored keys with frequency that depends on the data
density and on the requested $L$.
In-gap queries stay empty by
construction; near-key and uniform queries become non-empty when
their interval crosses a stored key. On Wikipedia timestamps with
$L=1024$ this gives a non-empty fraction of $\sim\!23\%$
(reflecting the dense $30$-bit universe); on Facebook user IDs
($33$-bit universe) the fraction is $\sim\!14\%$; on uniformly
spread keys at $n=2^{20}$ it is well below $1\%$.

The FPR estimator runs over the empty subset alone (the only
queries on which a false positive can be defined), while the
latency estimator times the whole stream and reports nanoseconds
per query across both empty and non-empty answers. The two query
sets share their RNG seed up to the truncation step, so the
FPR-side measurement is exactly the empty-only restriction of the
latency-side measurement; latency differences between the two
modes are therefore attributable to the filter handling
non-empty queries, not to a different sampling distribution.
